\uuid{qP0z}
\titre{Système linéaire/Système échelonné/Rang d'un système}
\auteur{hanani ; mzari}
\organisation{exo7}
\licence{CC BY-NC-SA 4.0}
\datecreate{2020-01-20}
\niveau{L1}
\module{Algèbre}
\chapitre{Déterminant, système linéaire}
\sousChapitre{Système linéaire, rang}
\difficulte{3}
\theme{Systèmes d'équations linéaires, Système linéaire, Système échelonné, Rang d'un système}
\exo7qcmid{lille-2.systemes.022}

\contenu{
\begin{question}
\qclassification{200.03}
\qtags{motcle=Système linéaire/Système échelonné/Rang d'un système}
Soit $(\mathtt{S})$ un système à $4$ équations et $3$ inconnues, $(\mathtt{S}_E)$ un système échelonné obtenu par la méthode de résolution du pivot de Gauss et $r$ le rang du système $(\mathtt{S})$, c.à.d le nombre d'équations de $(\mathtt{S}_E)$ ayant un premier membre non nul. Quelles sont les assertions vraies ?
\begin{answers}
\bad{Si $r=1$, alors $(\mathtt{S})$ admet une infinité de solutions.}
\good{Si $r=2$ et les équations ayant un premier membre nul admettent aussi le second membre nul, alors $(\mathtt{S})$ admet une infinité de solutions.}
\good{Si $r=3$ et l'équation ayant un premier membre nul admet aussi le second membre nul, alors $(\mathtt{S})$ admet une unique solution.}
\bad{Si $r=3$, alors $(\mathtt{S})$ admet une unique solution.}
\end{answers}
\begin{explanations} D'après la méthode du pivot de Gauss, $(\mathtt{S})$ admet une solution si et seulement si toute équation de $(\mathtt{S}_E)$ dont le premier membre est nul a aussi son second membre nul.
\end{explanations}
\end{question}
}
